\documentclass[11pt]{article}
\usepackage[utf8]{inputenc}
\usepackage[T1]{fontenc}
\usepackage{newtxtext,newtxmath}
\usepackage{amsmath}
\usepackage{multicol}
\usepackage{geometry}
\usepackage{enumitem}
\usepackage{xcolor}
\usepackage{titlesec}

% Configurações de layout
\geometry{a4paper, left=1cm, right=1cm, top=0.5cm, bottom=1.2cm}
\setlength{\columnseprule}{0.4pt}
\setlength{\baselineskip}{1.0\baselineskip}

% Cores para títulos
\titleformat{\section}{\normalfont\Large\color{blue}}{\thesection}{1em}{}
\titleformat{\subsection}{\normalfont\large\color{red}}{\thesubsection}{1em}{}

\title{\textcolor{blue}{Gabarito: \textit{Men Against Fire} (Black Mirror)}}
\author{Professor: Jefferson}
\date{}

\begin{document}

\maketitle
\vspace{-1cm}

\begin{multicols}{2}

\section*{Parte 1: Compreensão}

\textbf{1. Função da MÁSCARA:} Otimizar o desempenho militar (exatidão, tática) e, principalmente, alterar a percepção sensorial para desumanizar o inimigo.

\textbf{2. Visualização das baratas:} Monstros mutantes com aparência cadavérica, dentes afiados e sons animalescos.

\textbf{3. Falha no implante:} Um dispositivo de luz emitido por uma "barata" interfere no sinal do implante MASS de Stripe.

\textbf{4. Identidade real:} São seres humanos comuns (homens, mulheres e crianças) que possuem "deficiências" genéticas ou pertencem a grupos sociais marginalizados.

\textbf{5. Sentidos:} O sistema bloqueia o cheiro de sangue/morte (substituindo por cheiros agradáveis ou neutros) e distorce os gritos de socorro em sons de monstros.

\textbf{6. Justificativa:} Limpeza genética; evitar que "genes inferiores" ou doenças se espalhem pela população humana.

\textbf{7. Dificuldade de matar:} Arquette explica que, historicamente, a maioria dos soldados não atira para matar devido à empatia humana; a tecnologia remove essa barreira.

\textbf{8. População civil:} Tratam as baratas com ódio e desprezo, acreditando na propaganda governamental de que elas são uma ameaça à saúde pública.

\textbf{9. Dilema de Stripe:} Ter suas memórias de assassinato resetadas e continuar servindo (ou viver na ilusão), ou ser preso revivendo o trauma real de suas ações.

\textbf{10. Cena final:} Stripe "vê" uma casa perfeita e uma vida feliz, mas a realidade mostra que ele está em frente a uma casa em ruínas, totalmente iludido pelo sistema.

\section*{Parte 2: Análise Crítica}

\textbf{11. Desumanização:} É usada para que o soldado não sinta culpa. Ao transformar o "outro" em algo não-humano, o assassinato torna-se apenas uma "limpeza".

\textbf{12. Eugenia:} O sistema justifica o extermínio através da lógica de que certas vidas têm menos valor genético, remetendo a ideologias nazistas e higienistas.

\textbf{13. Empatia:} A empatia depende do reconhecimento do outro como semelhante. Ao distorcer visão, voz e cheiro, o sistema impede qualquer conexão biológica ou emocional.

\textbf{14. Paralelo atual:} *Fake news* e discursos de ódio em redes sociais criam "inimigos" imaginários, retirando a humanidade de oponentes políticos ou sociais.

\textbf{15. Resposta pessoal:} Espera-se que o aluno reflita sobre como a obediência cega e a tecnologia podem anestesiar o senso moral individual.

\section*{Parte 3: Conexão Pessoal}

\textbf{Respostas Pessoais (16-19):} O objetivo é que o aluno relacione o episódio com a realidade das redes sociais, a importância da ética e a resistência contra manipulações institucionais.

\end{multicols}

\end{document}
